% Part: first-order-logic
% Chapter: syntax-and-semantics
% Section: covered-structures

\documentclass[../../../include/open-logic-section]{subfiles}

\begin{document}

\olfileid{fol}{syn}{cov}

\olsection{প্রথম-ক্রমের ভাষার আচ্ছাদিত \printtoken{P}{structure}}

\begin{explain}
মনে করুন, কোনো পদে !!{variable}s না থাকলে তাকে \emph{বদ্ধ} বলে।
\end{explain}

\begin{defn}[বদ্ধ পদের !!^{value}]
$t$ যদি $\Lang L$ ভাষার একটি বদ্ধ পদ এবং $\Struct M$ যদি $\Lang L$-এর
একটি !!{structure} হয়, তাহলে \emph{!!{value}}~$\Value{t}{M}$ নিচের
নিয়মে সংজ্ঞায়িত:
\begin{enumerate}
\item $t$ যদি শুধু !!{constant}~$c$ হয়, তাহলে $\Value{c}{M} = \Assign{c}{M}$।
\item $t$ যদি $\Atom{f}{t_1, \ldots, t_n}$ রূপের হয়, তাহলে
  \[
  \Value{t}{M} = \Assign{f}{M}(\Value{t_1}{M}, \ldots,
  \Value{t_n}{M}).
  \]
\end{enumerate}
\end{defn}

\begin{defn}[আচ্ছাদিত !!{structure}]
কোনো !!{structure}-এর সংজ্ঞাক্ষেত্রের প্রত্যেক সদস্য যদি কোনো বদ্ধ
পদের !!{value} হয়, তাহলে গঠনটি \emph{আচ্ছাদিত}।
\end{defn}

\begin{ex}
ধরি~$\Lang L$ এমন ভাষা যাতে !!{constant}s $\Obj{zero}$,
$\Obj{one}$, $\Obj{two}$, \dots, দ্বি-স্থানীয় !!{predicate}~$<$ এবং
দ্বি-স্থানীয় !!{function}s $+$ ও $\times$ আছে। তাহলে $\Lang L$-এর
!!a{structure}~$\Struct M$-এর সংজ্ঞাক্ষেত্র $\Domain M = \{0, 1, 2,
\ldots\}$ এবং আরোপগুলি $\Assign{\Obj{zero}}{M} = 0$,
$\Assign{\Obj{one}}{M} = 1$, $\Assign{\Obj{two}}{M} = 2$, ইত্যাদি।
দ্বি-স্থানীয় সম্বন্ধসংকেত $<$-এর জন্য $\Assign{<}{M}$ হলো এমন সব
জোড়া $\tuple{c_1, c_2} \in \Domain{M}^2$-এর সেট, যাদের ক্ষেত্রে
$c_1$,~$c_2$-এর চেয়ে ছোট: যেমন, $\tuple{1, 3} \in
\Assign{<}{M}$, কিন্তু $\tuple{2, 2} \notin \Assign{<}{M}$।
দ্বি-স্থানীয় !!{function} $+$-এর জন্য $\Assign{+}{M}$-কে প্রচলিত
উপায়ে সংজ্ঞায়িত করি—যেমন, $\Assign{+}{M}(2,3)$-এর প্রতিচিত্র~$5$;
দ্বি-স্থানীয় !!{function}~$\times$-এর ক্ষেত্রেও একইভাবে। তাই
$\Obj{four}$-এর !!{value} শুধু~$4$, আর
$\times(\Obj{two}, +(\Obj{three},\Obj{zero}))$-এর !!{value}
(অথবা মধ্যলিখিত সংকেতরীতিতে $\Obj{two} \times
(\Obj{three} + \Obj{zero})$) হলো
\begin{multline*}
\Value{\times(\Obj{two}, +(\Obj{three},\Obj{zero}))}{M} =\\
\begin{aligned}
& =\Assign{\times}{M}(\Value{\Obj{two}}{M}, \Value{+(\Obj{three}, \Obj{zero})}{M})\\
& = \Assign{\times}{M}(\Value{\Obj{two}}{M}, \Assign{+}{M}(\Value{\Obj{three}}{M},
\Value{\Obj{zero}}{M})) \\
& = \Assign{\times}{M}(\Assign{\Obj{two}}{M}, \Assign{+}{M}(\Assign{\Obj{three}}{M},
\Assign{\Obj{zero}}{M})) \\
& = \Assign{\times}{M}(2, \Assign{+}{M}(3, 0)) \\
& = \Assign{\times}{M}(2, 3) \\
& = 6
\end{aligned}
\end{multline*}
\end{ex}

\begin{prob}
পাটিগণিতের মানক মডেল $\Struct N$ কি আচ্ছাদিত? ব্যাখ্যা করুন।
\end{prob}

\end{document}
